% Practice list C --- topics of Listas 5-6 (Kurlat chs. 9-11).
% \solucoesfalse -> blank list ; \solucoestrue -> answer key + rubric.
% Compile: pdflatex (2x); check: pdffonts (all Type 1). Numbers: practice-lists-check.py
\documentclass[11pt,a4paper]{article}
% Shared preamble of the practice lists (practice-L*.tex), adapted from lista-exam.tex.
% Each list sets \solucoesfalse / \solucoestrue and \listatema right after % Shared preamble of the practice lists (practice-L*.tex), adapted from lista-exam.tex.
% Each list sets \solucoesfalse / \solucoestrue and \listatema right after % Shared preamble of the practice lists (practice-L*.tex), adapted from lista-exam.tex.
% Each list sets \solucoesfalse / \solucoestrue and \listatema right after \input{practice-preamble}.
\usepackage[utf8]{inputenc}
\usepackage[T1]{fontenc}
\usepackage{lmodern}   % Type 1 vector fonts (replaces PK bitmaps)
\usepackage{cmap}      % ToUnicode CMap -> copyable text and accents
\usepackage[english]{babel}
\usepackage{amsmath,amssymb,amsfonts}
\usepackage[a4paper,margin=2.2cm]{geometry}
\usepackage{enumitem}
\usepackage{fancyhdr}
\usepackage{comment}
\usepackage{xcolor}
\usepackage{microtype}
\usepackage{booktabs}
\usepackage{tikz}

\newif\ifsolucoes
\newcommand{\disciplina}{Macroeconomics I --- MPE Insper}
\newcommand{\espacoresol}{4.5cm}   % height of the answer space in the blank list
% The sol environment must be decided after the list sets its toggle.
\AtBeginDocument{\ifsolucoes\includecomment{sol}\else\excludecomment{sol}\fi}

\pagestyle{fancy}\fancyhf{}
\lhead{\footnotesize\disciplina}
\rhead{\footnotesize Practice list\ifsolucoes{} --- Solutions and rubric\fi}
\cfoot{\footnotesize\thepage}
\renewcommand{\headrulewidth}{0.3pt}
\setlength{\headheight}{14pt}
\setlist{topsep=2pt,itemsep=2pt}
\newcounter{questao}

\newcommand{\bloco}[3]{\bigskip\noindent
  {\large\bfseries\color{blue!55!black}Block #1 --- #2}\hfill\textit{(#3)}\par
  \rule{\linewidth}{0.3pt}\smallskip}
\newcommand{\questao}[1]{\refstepcounter{questao}\medskip\noindent
  \textbf{Question \thequestao.}~#1\par}
\newenvironment{alternativas}{\begin{enumerate}[label=\Alph*.,leftmargin=2.2em,topsep=1pt]}{\end{enumerate}}
\newcommand{\gabarito}[1]{\ifsolucoes\par\noindent\textbf{Answer:}~#1\fi}
\newcommand{\resolespaco}{\ifsolucoes\else\par\textit{Solution:}\vspace{\espacoresol}\par\fi}
\newcommand{\passo}[2]{\par\smallskip\noindent\textbf{Step --- #1.}~#2}
\newcommand{\intuicao}[1]{\par\smallskip\noindent\textit{Economic intuition.}~#1}
\newcommand{\refbib}[1]{\par\smallskip\noindent{\footnotesize\color{black!60}\textbf{Ref:} #1}}
\newcommand{\rubrica}[1]{\par\smallskip\noindent\textbf{Grading rubric.}\begin{itemize}[leftmargin=1.4em,topsep=1pt]#1\end{itemize}}

\newcommand{\cabecalho}{%
\begin{center}
  {\Large\bfseries \disciplina}\\[2pt]
  {\large Practice list: \listatema}\ifsolucoes\\[2pt]\textit{Solutions and grading rubric}\fi
\end{center}
\vspace{2pt}
{\footnotesize\noindent\textbf{Instructions.} Multiple choice: mark the single correct option (no
penalty for wrong answers). \emph{True/False}: every 2 wrong items cancel 1 right one
(score $=\max\{0,\text{right}-\text{wrong}/2\}\times$ item value); blank items do not count.
Discursive questions: show every step. Numbers are illustrative and chosen to give clean answers.\par}}
.
\usepackage[utf8]{inputenc}
\usepackage[T1]{fontenc}
\usepackage{lmodern}   % Type 1 vector fonts (replaces PK bitmaps)
\usepackage{cmap}      % ToUnicode CMap -> copyable text and accents
\usepackage[english]{babel}
\usepackage{amsmath,amssymb,amsfonts}
\usepackage[a4paper,margin=2.2cm]{geometry}
\usepackage{enumitem}
\usepackage{fancyhdr}
\usepackage{comment}
\usepackage{xcolor}
\usepackage{microtype}
\usepackage{booktabs}
\usepackage{tikz}

\newif\ifsolucoes
\newcommand{\disciplina}{Macroeconomics I --- MPE Insper}
\newcommand{\espacoresol}{4.5cm}   % height of the answer space in the blank list
% The sol environment must be decided after the list sets its toggle.
\AtBeginDocument{\ifsolucoes\includecomment{sol}\else\excludecomment{sol}\fi}

\pagestyle{fancy}\fancyhf{}
\lhead{\footnotesize\disciplina}
\rhead{\footnotesize Practice list\ifsolucoes{} --- Solutions and rubric\fi}
\cfoot{\footnotesize\thepage}
\renewcommand{\headrulewidth}{0.3pt}
\setlength{\headheight}{14pt}
\setlist{topsep=2pt,itemsep=2pt}
\newcounter{questao}

\newcommand{\bloco}[3]{\bigskip\noindent
  {\large\bfseries\color{blue!55!black}Block #1 --- #2}\hfill\textit{(#3)}\par
  \rule{\linewidth}{0.3pt}\smallskip}
\newcommand{\questao}[1]{\refstepcounter{questao}\medskip\noindent
  \textbf{Question \thequestao.}~#1\par}
\newenvironment{alternativas}{\begin{enumerate}[label=\Alph*.,leftmargin=2.2em,topsep=1pt]}{\end{enumerate}}
\newcommand{\gabarito}[1]{\ifsolucoes\par\noindent\textbf{Answer:}~#1\fi}
\newcommand{\resolespaco}{\ifsolucoes\else\par\textit{Solution:}\vspace{\espacoresol}\par\fi}
\newcommand{\passo}[2]{\par\smallskip\noindent\textbf{Step --- #1.}~#2}
\newcommand{\intuicao}[1]{\par\smallskip\noindent\textit{Economic intuition.}~#1}
\newcommand{\refbib}[1]{\par\smallskip\noindent{\footnotesize\color{black!60}\textbf{Ref:} #1}}
\newcommand{\rubrica}[1]{\par\smallskip\noindent\textbf{Grading rubric.}\begin{itemize}[leftmargin=1.4em,topsep=1pt]#1\end{itemize}}

\newcommand{\cabecalho}{%
\begin{center}
  {\Large\bfseries \disciplina}\\[2pt]
  {\large Practice list: \listatema}\ifsolucoes\\[2pt]\textit{Solutions and grading rubric}\fi
\end{center}
\vspace{2pt}
{\footnotesize\noindent\textbf{Instructions.} Multiple choice: mark the single correct option (no
penalty for wrong answers). \emph{True/False}: every 2 wrong items cancel 1 right one
(score $=\max\{0,\text{right}-\text{wrong}/2\}\times$ item value); blank items do not count.
Discursive questions: show every step. Numbers are illustrative and chosen to give clean answers.\par}}
.
\usepackage[utf8]{inputenc}
\usepackage[T1]{fontenc}
\usepackage{lmodern}   % Type 1 vector fonts (replaces PK bitmaps)
\usepackage{cmap}      % ToUnicode CMap -> copyable text and accents
\usepackage[english]{babel}
\usepackage{amsmath,amssymb,amsfonts}
\usepackage[a4paper,margin=2.2cm]{geometry}
\usepackage{enumitem}
\usepackage{fancyhdr}
\usepackage{comment}
\usepackage{xcolor}
\usepackage{microtype}
\usepackage{booktabs}
\usepackage{tikz}

\newif\ifsolucoes
\newcommand{\disciplina}{Macroeconomics I --- MPE Insper}
\newcommand{\espacoresol}{4.5cm}   % height of the answer space in the blank list
% The sol environment must be decided after the list sets its toggle.
\AtBeginDocument{\ifsolucoes\includecomment{sol}\else\excludecomment{sol}\fi}

\pagestyle{fancy}\fancyhf{}
\lhead{\footnotesize\disciplina}
\rhead{\footnotesize Practice list\ifsolucoes{} --- Solutions and rubric\fi}
\cfoot{\footnotesize\thepage}
\renewcommand{\headrulewidth}{0.3pt}
\setlength{\headheight}{14pt}
\setlist{topsep=2pt,itemsep=2pt}
\newcounter{questao}

\newcommand{\bloco}[3]{\bigskip\noindent
  {\large\bfseries\color{blue!55!black}Block #1 --- #2}\hfill\textit{(#3)}\par
  \rule{\linewidth}{0.3pt}\smallskip}
\newcommand{\questao}[1]{\refstepcounter{questao}\medskip\noindent
  \textbf{Question \thequestao.}~#1\par}
\newenvironment{alternativas}{\begin{enumerate}[label=\Alph*.,leftmargin=2.2em,topsep=1pt]}{\end{enumerate}}
\newcommand{\gabarito}[1]{\ifsolucoes\par\noindent\textbf{Answer:}~#1\fi}
\newcommand{\resolespaco}{\ifsolucoes\else\par\textit{Solution:}\vspace{\espacoresol}\par\fi}
\newcommand{\passo}[2]{\par\smallskip\noindent\textbf{Step --- #1.}~#2}
\newcommand{\intuicao}[1]{\par\smallskip\noindent\textit{Economic intuition.}~#1}
\newcommand{\refbib}[1]{\par\smallskip\noindent{\footnotesize\color{black!60}\textbf{Ref:} #1}}
\newcommand{\rubrica}[1]{\par\smallskip\noindent\textbf{Grading rubric.}\begin{itemize}[leftmargin=1.4em,topsep=1pt]#1\end{itemize}}

\newcommand{\cabecalho}{%
\begin{center}
  {\Large\bfseries \disciplina}\\[2pt]
  {\large Practice list: \listatema}\ifsolucoes\\[2pt]\textit{Solutions and grading rubric}\fi
\end{center}
\vspace{2pt}
{\footnotesize\noindent\textbf{Instructions.} Multiple choice: mark the single correct option (no
penalty for wrong answers). \emph{True/False}: every 2 wrong items cancel 1 right one
(score $=\max\{0,\text{right}-\text{wrong}/2\}\times$ item value); blank items do not count.
Discursive questions: show every step. Numbers are illustrative and chosen to give clean answers.\par}}

\solucoesfalse                 % <<<< \solucoestrue for the ANSWER KEY
\newcommand{\listatema}{General Equilibrium, Money and Inflation (Listas 5--6)}

\begin{document}
\cabecalho

% =============================================================================
\bloco{1}{General equilibrium}{Kurlat ch.~9 \,\textbullet\, 25 points}

\questao{(10 points) Objective items. The Tesouro asks what pins down the real rate on NTN-Bs in
model economies with a representative household and competitive markets.}

\textbf{a)} (3 points) Two-period endowment economy, $u=\ln c$, $\beta=0.96$, $y_1=100$, $y_2=104$,
no government. The equilibrium real interest rate is:
\begin{alternativas}
\item $4.00\%$
\item $8.33\%$
\item $4.17\%$
\item $0.00\%$
\end{alternativas}
\gabarito{B}
\begin{sol}
\passo{market clearing}{The representative household cannot trade with itself: $a=0$,
$c_t=y_t$. Euler: $1+r=\dfrac{u'(y_1)}{\beta u'(y_2)}=\dfrac{y_2}{\beta y_1}=\dfrac{104}{96}=1.0833$.
Option C is $1/\beta-1$, which ignores growth.}
\refbib{Kurlat (2020), \S9.1, p.~165.}
\end{sol}

\medskip
\textbf{b)} (3 points) Infinite-horizon economy on a balanced growth path where consumption
grows at $g=2\%$, $u=c^{1-\sigma}/(1-\sigma)$, $\sigma=2$, $\beta=0.98$. The real interest rate is:
\begin{alternativas}
\item $2.04\%$
\item $4.08\%$
\item $8.24\%$
\item $6.16\%$
\end{alternativas}
\gabarito{D}
\begin{sol}
\passo{Euler}{$(1+g)^{\sigma}=\beta(1+r)\Rightarrow1+r=1.02^2/0.98=1.0616$. Approximation:
$r\approx\rho+\sigma g=2.04+4=6.04\%$.}
\refbib{Kurlat (2020), \S9.3, p.~172.}
\end{sol}

\medskip
\textbf{c)} (2 points, T or F) In the two-period endowment economy of (a), a rise in patience
($\beta\uparrow$) raises equilibrium saving.
\gabarito{F}
\begin{sol}
\passo{price adjusts}{Aggregate saving is zero in any equilibrium. Higher $\beta$ raises the
desire to save at the old $r$, so $r$ falls until households again want $a=0$.}
\refbib{Kurlat (2020), \S9.1.}
\end{sol}

\medskip
\textbf{d)} (2 points, T or F) By the First Welfare Theorem, the steady-state capital stock of the
competitive infinite-horizon economy maximizes steady-state consumption (Golden Rule).
\gabarito{F}
\begin{sol}
\passo{modified rule}{The competitive steady state satisfies $F_K-\delta=1/\beta-1$, below the
Golden Rule's $F_K-\delta=0$. It is Pareto-efficient because impatient households value
consumption now; maximizing \emph{steady-state} consumption is not their objective.}
\refbib{Kurlat (2020), \S9.2--9.3.}
\end{sol}

% -----------------------------------------------------------------------------
\questao{(15 points) Discursive. Infinite-horizon economy: household with
$\sum_t\beta^t c_t^{1-\sigma}/(1-\sigma)$, inelastic labor $L_t=1$; firm with $Y=AK^{\alpha}L^{1-\alpha}$;
capital depreciates at $\delta$; $K_0$ given. Numbers: $\beta=0.95$, $\sigma=2$, $\alpha=1/3$,
$\delta=0.1$, $A=1$.}

\textbf{a)} (5 points) Write the household's period budget constraint when it owns and rents the
capital, and derive its Euler equation. Solve the firm's problem. Show that $1+r_{t+1}=r^K_{t+1}+1-\delta$
and define a competitive equilibrium (allocation, prices, conditions).
\resolespaco
\begin{sol}
\passo{household}{$c_t+K_{t+1}=w_t+r^K_tK_t+(1-\delta)K_t$ (plus bonds $b_t$ in zero net supply:
$+\,(1+r_t)b_t-b_{t+1}$).}
\passo{Euler}{$c_t^{-\sigma}=\beta(1+r_{t+1})c_{t+1}^{-\sigma}$. The bond and capital FOCs give the
same Euler only if $1+r_{t+1}=r^K_{t+1}+1-\delta$ (no arbitrage).}
\passo{firm}{$\max AK^{\alpha}L^{1-\alpha}-w L-r^K K$: $r^K=\alpha AK^{\alpha-1}$, $w=(1-\alpha)AK^{\alpha}$;
CRS gives zero profits.}
\passo{definition}{Allocation $\{c_t,K_{t+1},L_t\}$ and prices $\{w_t,r^K_t,r_t\}$ such that
(i) the household maximizes utility given prices; (ii) the firm maximizes profits given prices;
(iii) markets clear: $L_t=1$, capital demanded $=K_t$, bonds $b_t=0$, goods
$c_t+K_{t+1}-(1-\delta)K_t=AK_t^{\alpha}$.}
\intuicao{Bonds and capital are two ways to carry resources forward; in equilibrium they must
pay the same, or no one would hold one of them.}
\refbib{Kurlat (2020), \S9.3, pp.~172--180; Lista 5, Q2.}
\rubrica{
\item Budget constraint and Euler: 1.5 pts.
\item Firm FOCs: 1 pt.
\item No-arbitrage derived, not just stated: 1 pt.
\item Definition with all three conditions and every market: 1.5 pts. Missing market clearing:
  max.\ 3 pts in the item.
}
\end{sol}

\textbf{b)} (5 points) Compute the steady state: $r$, $r^K$, $K$, $Y$, $C$, $w$ and the saving rate
$\delta K/Y$. Why does $\sigma$ not appear, and why does $r$ depend only on $\beta$?
\resolespaco
\begin{sol}
\passo{interest}{$c_{t+1}=c_t\Rightarrow1+r=1/\beta$: $r=5.26\%$; $r^K=r+\delta=0.1526$.}
\passo{capital}{$\alpha K^{\alpha-1}=0.1526\Rightarrow K=(0.3333/0.1526)^{3/2}=3.227$.}
\passo{rest}{$Y=K^{1/3}=1.478$; $C=Y-\delta K=1.155$; $w=\tfrac23Y=0.985$;
$\delta K/Y=21.8\%$.}
\passo{$\sigma$}{$\sigma$ governs the willingness to deviate from flat consumption. In steady state
consumption is flat, so only $\beta$ matters; $\sigma$ shapes the transition, not the destination.}
\intuicao{A household that keeps consumption constant needs a return that exactly compensates
its impatience; technology then adjusts $K$ until the marginal product delivers it.}
\refbib{Kurlat (2020), \S9.3.}
\rubrica{
\item $r$ and $r^K$: 1.5 pts.
\item $K$, $Y$, $C$, $w$, saving rate: 2.5 pts. Forgetting $\delta$ in $r^K$: max.\ 2 pts in the item.
\item Role of $\sigma$ and $\beta$: 1 pt.
}
\end{sol}

\textbf{c)} (5 points) Compute the Golden Rule $K_{gr}$ and $C_{gr}$ and show $K<K_{gr}$. Explain
why this is consistent with the First Welfare Theorem and what happens as $\beta\to1$. Then let
$A$ rise permanently to $1.1$: by how much do $K$, $Y$, $C$, $w$ and $r$ change in the new steady state?
\resolespaco
\begin{sol}
\passo{Golden Rule}{$\max_K K^{\alpha}-\delta K\Rightarrow\alpha K^{\alpha-1}=\delta$:
$K_{gr}=(0.3333/0.1)^{3/2}=6.086$, $C_{gr}=1.217>1.155$.}
\passo{why lower}{$F_K=r+\delta>\delta$ and $F_K$ is decreasing, so $K<K_{gr}$. Getting to $K_{gr}$
would cost consumption now; impatient households do not want that trade. As $\beta\to1$,
$r\to0$ and $K\to K_{gr}$.}
\passo{TFP}{$r$ is still $1/\beta-1$, so $r^K$ is unchanged and $K/Y$ too.
$K=(\alpha A/r^K)^{1/(1-\alpha)}$ scales with $A^{3/2}$: $K$, $Y$, $C$ and $w$ all rise by
$1.1^{1.5}-1=15.4\%$; $r$ unchanged.}
\intuicao{The planner and the market agree: both discount the future at $\beta$. The Golden Rule
ignores the time cost of building capital.}
\refbib{Kurlat (2020), \S9.2 and \S9.3; \S4.3 (Golden Rule).}
\rubrica{
\item $K_{gr}$, $C_{gr}$ and the comparison: 1.5 pts.
\item Consistency with the FWT (impatience) and $\beta\to1$: 2 pts. Calling the economy
  inefficient: max.\ 2 pts in the item.
\item TFP comparative statics with $r$ unchanged: 1.5 pts.
}
\end{sol}

% =============================================================================
\bloco{2}{Money and inflation}{Kurlat chs.~10--11 \,\textbullet\, 25 points}

\questao{(10 points) Objective items. Banco Central data and money-demand theory.}

\textbf{a)} (3 points) The public holds currency equal to 25\% of demand deposits and banks
hold reserves equal to 15\% of deposits. With a monetary base of R\$ 200 billion, M1 is:
\begin{alternativas}
\item R\$ 500 billion
\item R\$ 1{,}333 billion
\item R\$ 625 billion
\item R\$ 250 billion
\end{alternativas}
\gabarito{C}
\begin{sol}
\passo{multiplier}{$m=\dfrac{1+cr}{rr+cr}=\dfrac{1.25}{0.40}=3.125$; $M1=3.125\cdot200=625$.
Option B ($200/0.15$) ignores currency leakage.}
\refbib{Kurlat (2020), \S10.2--10.3, pp.~192--199.}
\end{sol}

\medskip
\textbf{b)} (3 points) Real money demand is $m=e^{-a\pi}$ with $a=4$ (Cagan). The steady-state
inflation that maximizes revenue from the inflation tax $\pi m$ is:
\begin{alternativas}
\item $25\%$
\item $40\%$
\item $4\%$
\item $100\%$
\end{alternativas}
\gabarito{A}
\begin{sol}
\passo{Laffer}{$\frac{d}{d\pi}\pi e^{-a\pi}=e^{-a\pi}(1-a\pi)=0\Rightarrow\pi=1/a=25\%$.}
\refbib{Kurlat (2020), \S11.3, p.~215; Exercise 11.6.}
\end{sol}

\medskip
\textbf{c)} (2 points, T or F) If money is superneutral, permanently raising money growth from
4\% to 8\% leaves the nominal interest rate unchanged.
\gabarito{F}
\begin{sol}
\passo{Fisher}{Superneutrality concerns \emph{real} variables: $r$ is unchanged, so inflation
rises 4 points and $i=r+\pi$ rises 4 points too (Fisher effect).}
\refbib{Kurlat (2020), \S11.2, p.~208.}
\end{sol}

\medskip
\textbf{d)} (2 points, T or F) With Baumol--Tobin demand $M/P=\sqrt{FY/(2i)}$ ($F$ a real cost
per trip), if real income and the nominal interest rate both double, real balances are unchanged.
\gabarito{T}
\begin{sol}
\passo{elasticities}{Income elasticity $+\tfrac12$, interest elasticity $-\tfrac12$: the two
doublings cancel inside the square root.}
\refbib{Kurlat (2020), \S10.4, p.~199.}
\end{sol}

% -----------------------------------------------------------------------------
\questao{(15 points) Discursive. Brazil in the 1980s. Real money demand is
$M/P=\kappa\,Y e^{-a i}$ with $i=r+\pi$; $Y$ grows at $g$, money at $\mu$. Numbers: $\kappa=0.3$,
$a=2.5$, $r=0.04$, $g=0.02$ (illustrative).}

\textbf{a)} (5 points) Show that in a steady state with constant $i$, $\pi=\mu-g$. Write
seigniorage $\dot M/P$ and the inflation tax $\pi M/P$ as shares of $Y$ as functions of $\mu$, and
explain their difference. When does the approximation $i=r+\pi$ fail?
\resolespaco
\begin{sol}
\passo{inflation}{Constant $i$ makes $M/(PY)$ constant, so $\mu-\pi-g=0$: $\pi=\mu-g$.}
\passo{seigniorage}{$\dfrac{\dot M}{PY}=\mu\dfrac{M}{PY}=\mu\kappa e^{-a(r+\mu-g)}$;
inflation tax $=\pi\dfrac{M}{PY}=(\mu-g)\kappa e^{-a(r+\mu-g)}$.}
\passo{difference}{$g\,M/(PY)$: a growing economy demands more real balances, which the
government can issue without inflation. This part is not a tax on holders.}
\passo{Fisher}{Exact: $1+i=(1+r)(1+\pi)$; the cross term $r\pi$ is negligible only at low
inflation. At 80\% a year it is not.}
\refbib{Kurlat (2020), \S11.2--11.3, pp.~208--217.}
\rubrica{
\item Derives $\pi=\mu-g$ from constant real balances per unit of output: 1.5 pts.
\item Both revenue formulas: 1.5 pts. Treating seigniorage and inflation tax as two additive
  revenues: max.\ 2 pts in the item.
\item Difference and Fisher caveat: 2 pts.
}
\end{sol}

\textbf{b)} (5 points) Find the money growth rate that maximizes seigniorage, the associated
inflation, the maximum revenue (\% of GDP) and the share collected as inflation tax.
\resolespaco
\begin{sol}
\passo{FOC}{$\dfrac{d}{d\mu}\mu e^{-a\mu}=e^{-a\mu}(1-a\mu)=0\Rightarrow\mu^*=1/a=40\%$.}
\passo{numbers}{$\pi=38\%$; $S/Y=0.4\cdot0.3\cdot e^{-2.5(0.04+0.38)}=0.12\,e^{-1.05}=4.20\%$ of GDP.
Inflation-tax share $=\pi/\mu=95\%$.}
\intuicao{Above $\mu^*$ the base $M/P$ shrinks faster than the rate rises: the Laffer curve of
inflation. Money printing has a ceiling on what it can finance.}
\refbib{Kurlat (2020), \S11.3; Exercise 11.6 (Cagan).}
\rubrica{
\item FOC and $\mu^*=1/a$: 2 pts.
\item Inflation, revenue and share: 2 pts.
\item Laffer interpretation: 1 pt.
}
\end{sol}

\textbf{c)} (5 points) The government must finance a deficit of 3\% of GDP by printing money.
Find the steady states numerically and say which one is on the wrong side of the Laffer
curve. What inflation does the Friedman rule prescribe here, and why do central banks such as
the BCB target positive inflation instead?
\resolespaco
\begin{sol}
\passo{two roots}{Solve $\mu\cdot0.3\,e^{-2.5(0.02+\mu)}=0.03$: $\mu=15.5\%$ ($\pi=13.5\%$) and
$\mu=82.3\%$ ($\pi=80.3\%$). Both finance the deficit; the second is beyond $\mu^*=40\%$.}
\passo{wrong side}{At $\pi=80\%$ the same revenue is raised with a much smaller base and much
larger distortions (shoe-leather, menu, relative-price dispersion): strictly worse.}
\passo{Friedman rule}{Set $i=0$: $\pi=-r=-4\%$, removing the opportunity cost of money.}
\passo{why positive targets}{Measurement bias in price indices, downward nominal rigidity, and
room to cut $i$ before hitting zero.}
\intuicao{A deficit that money must finance pins the inflation rate; the only choice left is
which side of the Laffer curve, and expectations can push the economy to the bad one.}
\refbib{Kurlat (2020), \S11.3--11.4, pp.~215--222.}
\rubrica{
\item Both roots: 2 pts. Only one root found: max.\ 3 pts in the item.
\item Identifies the bad equilibrium and why: 1 pt.
\item Friedman rule: 1 pt.
\item Reasons for positive targets: 1 pt.
}
\end{sol}

\end{document}
